% =============================================================================
% D5 — Demande, préférences et élasticité — v2 (24 mai 2026)
% Ce fichier contient la D5 complète à coller dans Chapter02.tex
% en remplacement de tout le contenu entre
%   \subsection{Demande, préférences et élasticité}
% et
%   \subsection{Signal de prix et structure des marchés}
% =============================================================================
% ATTENTION : Chapter02.tex a été partiellement édité par Claude.
% Faire d'abord : git checkout -- Chapters/Chapter02.tex
% Puis remplacer l'ancienne D5 par le contenu ci-dessous.
% Les entrées bib manquantes sont dans d5_bib_entries.bib (à merger).
% =============================================================================

% --- COPIER À PARTIR D'ICI ---

\subsection{Demande, préférences et élasticité}\label{sec:d5-demande}

Le cadre standard de la demande suppose que le consommateur observe le coût de ce qu'il consomme, directement ou par le prix. Les marchés de biens informationnels peuvent violer cette hypothèse de trois manières~: le prix perçu peut être nul alors que le coût de production est positif, l'usage peut consommer du temps plutôt que du budget, et la demande elle-même peut être produite par le système qui la sert. Cette dimension examine comment ces propriétés transforment la structure de la demande finale et les signaux auxquels les agents économiques répondent.

\paragraph{Contributions propres de la section}
Cette section identifie trois mécanismes économiquement distincts par lesquels l'ICT modifie la structure de la demande finale. Le premier~--- H-D5.1~--- est la dissociation entre le prix perçu et le coût de production, qui rompt le bouclage prix-rareté du cadre standard et crée une asymétrie d'élasticité entre utilisateurs finaux et opérateurs d'infrastructure. Le second~--- H-D5.2~--- est la réallocation du temps comme contrainte active, canal absent du cadre où le temps n'est pas une contrainte de consommation~; ce canal mute avec l'IA agentique, qui déplace la borne du temps humain vers le budget de calcul. Le troisième~--- H-D5.3~--- est la production de la demande par l'architecture informationnelle, qui distingue la demande induite de la demande révélée et dont les conséquences en compute sont directes. Les trois opèrent conjointement et produisent un résultat que ni l'un ni l'autre ne produit seul.


\subsubsection{Le champ~: la demande au-delà des élasticités-prix}

Le cadre standard de la demande~--- préférences stables, contrainte de budget, prix comme signal de rareté~--- suppose que le consommateur observe le coût de ce qu'il consomme, directement ou par le prix. La demande est l'expression de préférences préexistantes, révélées par les choix observés sous contrainte de budget. L'information transforme cette relation par trois canaux que la littérature traite séparément mais qui opèrent conjointement dans les marchés informationnels contemporains.

Le premier canal est celui de la dissociation entre prix et coût. La littérature des biens d'expérience et des marchés bifaces, depuis \textcite{shapiro_information_1999} et \textcite{rochet_platform_2003}, a établi que les marchés de biens informationnels adoptent des structures de prix non standard~: un côté est subventionné, l'autre paie. \textcite{brynjolfsson_using_2019} ont quantifié le surplus associé à cette structure par des expériences de choix massives. Le canal est identifié~; ce qui manque est l'analyse des conséquences sur la transmission du signal de rareté et sur l'incidence du coût.

Le deuxième canal est celui du temps. \textcite{goolsbee_valuing_2006} proposent de valoriser les produits numériques par le temps passé plutôt que par la dépense. \textcite{rachel_leisure_2024} formalise un changement technique biaisé vers le loisir dans un modèle macroéconomique où les services numériques gratuits absorbent du temps en échange de l'attention. Le canal est formalisé~; ce qui manque est l'analyse de sa mutation quand la borne du temps humain est relâchée par l'automatisation de la demande.

Le troisième canal est celui de la formation de la demande par le système qui la sert. \textcite{anderson_competition_2012} modélisent l'attention comme une ressource commune sur-exploitée par les plateformes~; \textcite{allcott_digital_2022} quantifient la part de l'usage des réseaux sociaux attribuable à des problèmes de self-control (31~\%). Le canal est récent dans la littérature~; il n'a pas été articulé avec les deux précédents ni avec les observations historiques du broadcasting et de la mécanographie.


\begin{table}[!ht]
\caption{Trois canaux de transformation de la demande et leur variation entre configurations.}\label{tab:d5-canaux}
\centering\footnotesize
\renewcommand{\arraystretch}{1.35}
\begin{tabularx}{\textwidth}{@{}p{2.2cm}Xp{3.2cm}p{3.2cm}@{}}
\toprule
\textbf{Canal} & \textbf{Question simple} & \textbf{Où opère la rupture avec le cadre standard} & \textbf{Cas de référence} \\
\midrule
Dissociation prix-coût & Le consommateur observe-t-il le coût de ce qu'il consomme~? & Bouclage prix-rareté~: le prix ne transmet plus le signal de rareté au consommateur final & Broadcasting radio (1920s), Google AdWords (2001), forfait illimité mobile \\[4pt]
Réallocation du temps & L'usage consomme-t-il du budget ou du temps~? & Contrainte active~: quand le prix est nul, la contrainte de budget ne mord plus~; le temps disponible devient la borne & Télévision (Gentzkow 2006), réseaux sociaux (Allcott 2020), IA agentique (post-2024) \\[4pt]
Demande induite & La demande préexiste-t-elle au service, ou le service la produit-il~? & Nature des préférences~: la demande est endogène à l'architecture du système & IBM mécanographie (Bolle 1930s), AdSense (2003), autoplay, recommandation algorithmique \\
\bottomrule
\end{tabularx}
\renewcommand{\arraystretch}{1.0}
\end{table}


\subsubsection{Ce que le matériau révèle}

Avant de formuler les mécanismes, il est utile de poser ce que le matériau de l'analyse historique fait apparaître quand on le relit à travers la grille du tableau~\ref{tab:d5-canaux}.

Un premier constat concerne l'ancienneté de la dissociation prix-coût. Le broadcasting commercial américain des années~1920 invente le modèle~: la radio est gratuite pour l'auditeur, financée par la publicité~; le consommateur reçoit un service dont le coût de production est réel~--- émetteurs, studios, réseau de transmission~--- sans qu'aucun prix ne lui soit demandé. \textcite{braithwaite_economic_1928} analyse les effets économiques de la publicité dans l'\emph{Economic Journal} au moment précis où ce modèle s'installe. NBC (1926) et CBS (1928) organisent le broadcasting en réseaux nationaux financés par la publicité~: la structure de revenus tiers est en place un siècle avant les plateformes numériques. Plus ancien encore, la Navy américaine finance le réseau de stations du Pacifique pour 1,5~million de dollars en~1911 (section~\ref{sec:wireless})~: l'utilisateur militaire ne paie pas le coût marginal, l'État absorbe le fixe. Le même pattern se retrouve dans le Census pour Hollerith (section~\ref{sec:mechanography}) et la Convention pour le télégraphe Chappe (section~\ref{sec:telegraph})~: l'État crée le marché avant que la demande commerciale n'existe. La dissociation prix-coût n'est donc pas une invention du numérique~; elle est un trait récurrent des marchés informationnels, qui prend des formes différentes selon que le financeur tiers est l'État, l'annonceur ou l'opérateur de données.

Un deuxième constat concerne la capture du temps. \textcite{gentzkow_television_2006}, exploitant la variation géographique de l'introduction de la télévision aux États-Unis entre 1940 et 1960, montre que la télévision réduit la participation électorale d'environ 2~points de pourcentage par décennie, par un mécanisme de \emph{crowding out} de l'information politique~: la télévision, en tant que bien de divertissement supérieur, détourne le temps des journaux et de la radio. \textcite{stromberg_radio_2004} obtient le résultat inverse pour la radio~: celle-ci augmente la participation, parce qu'elle enrichit l'information politique sans détourner autant de la lecture. Le même modèle économique~--- gratuit, financé par la publicité~--- produit des effets opposés sur l'allocation du temps selon l'architecture du médium. Ce contraste est un test empirique du deuxième canal~: la réallocation du temps n'est pas une conséquence mécanique du prix nul~; elle dépend de la structure de substitution entre les services.

Un troisième constat concerne la production de la demande par le fournisseur. IBM, dans la période mécanographique, ne répond pas à une demande de calcul existante. Le bureau d'études de Bull au PLM~--- et l'équivalent chez IBM aux États-Unis~--- formule la devise «~d'abord organiser, ensuite mécaniser~» (section~\ref{sec:mechanography})~: le fournisseur reformate les processus du client, crée le besoin, puis y répond par son équipement. Les cartes perforées sont fournies gratuitement à l'entrée~; les données s'accumulent dans un format propriétaire qui rend la migration prohibitive. La demande de mécanographie est produite par l'infrastructure que le fournisseur déploie, pas par une préférence préexistante du client. Ce pattern se retrouve dans la séquence 2001--2003 documentée dans la section~\ref{subsec:turning_point}~: Google découvre que les traces comportementales~--- clics, requêtes, durées de session~--- sont des actifs prédictifs valorisables sur un marché publicitaire qui se constitue par leur agrégation. AdWords (2000, repensé en 2002) puis AdSense (2003) transforment l'attention des utilisateurs en matière première d'un nouveau type de marché. Le service est conçu pour maximiser le temps passé, pas pour satisfaire une préférence~: c'est de la demande induite par architecture.

Ces trois constats sont des régularités, pas encore des mécanismes. Les formuler comme mécanismes exige de préciser les conditions sous lesquelles ils opèrent et les cadres théoriques qui en rendent compte.


\subsubsection{Trois mécanismes causaux}

\paragraph{La dissociation prix-coût et l'asymétrie d'élasticité}

\markNEO \markINFO

Le premier canal est la dissociation entre le prix perçu par le consommateur et le coût de production du service. Le phénomène est plus ancien qu'on ne le croit. Le broadcasting commercial américain des années~1920 invente le modèle~; ce qui est nouveau dans l'économie numérique n'est pas le principe mais l'ampleur de la dissociation et la précision du mécanisme de financement tiers. \textcite{brynjolfsson_using_2019} estiment par des expériences de choix massives que le surplus consommateur de la recherche en ligne, de l'email et des réseaux sociaux représente des centaines de milliards de dollars\footnote{L'estimation est obtenue par \emph{willingness-to-accept} (WTA), méthodologie qui produit des valeurs typiquement plus élevées que la \emph{willingness-to-pay} (WTP). \textcite{allcott_welfare_2020} obtiennent des valorisations plus modestes pour les réseaux sociaux par d'autres protocoles. L'ordre de grandeur (surplus non négligeable, invisible dans les comptes nationaux) n'est pas contesté~; sa valeur précise l'est.}~--- un surplus qui n'apparaît dans aucune statistique de production parce qu'il est associé à un prix nul.

La dissociation a une conséquence sur la transmission du signal de rareté que la littérature sur les marchés bifaces ne développe pas entièrement. \textcite{bergemann_bonatti_2024} montrent que la plateforme monopolistique qui utilise les données pour orienter la demande extrait le surplus des deux côtés du marché~--- vendeurs et acheteurs~--- grâce à son avantage informationnel. Le coût de production du service (compute, stockage, bande passante, énergie) n'est pas supporté par l'utilisateur final~; il est absorbé par les marges de l'opérateur, transféré aux annonceurs qui paient pour l'accès, ou socialisé via l'infrastructure publique (subventions fiscales aux data centers, tarifs énergétiques préférentiels). L'incidence du coût~--- qui paie effectivement le compute~--- est une question de distribution dont les conséquences macroéconomiques diffèrent selon le canal~: si le coût est absorbé par les marges des hyperscalers, l'effet sur la demande finale est différent de celui qu'il aurait si le coût était transféré aux entreprises utilisatrices ou aux consommateurs.

La dissociation produit une asymétrie d'élasticité. \textcite{aubouin_economie_2023}, sur données INSEE 2007--2019, montre que les ménages au forfait illimité mobile ne modulent plus leur consommation physique avec le prix~: l'élasticité-prix de la consommation tombe. Le forfait illimité est un mode d'appropriation qui supprime le signal prix marginal~; une fois que le ménage ou la firme est au forfait, l'élasticité-prix s'effondre non seulement à la baisse mais aussi à la hausse, parce que le coût de migration vers un autre fournisseur excède le bénéfice d'un prix marginal plus bas. \textcite{axenbeck_what_2022} documentent un résultat partiellement symétrique côté firmes~: les entreprises plus numérisées sont en moyenne moins sensibles au prix de l'énergie dans leurs décisions d'investissement, parce que la structure tarifaire~--- contrats long terme, tarifs de gros, PPAs~--- lisse ou masque le signal. L'asymétrie (élastique à la baisse du prix de l'ICT, inélastique à la hausse du prix de l'énergie à cause du lock-in) est un trait du design du capital au sens de H-D2.2, pas seulement du prix nul.

\paragraph{Hypothèse H-D5.1} \emph{La dissociation entre le prix perçu par le consommateur et le coût de production rompt le bouclage prix-rareté qui structure le cadre standard de la demande, et crée une asymétrie d'élasticité entre utilisateurs finaux (faiblement élastiques au prix de l'énergie) et opérateurs d'infrastructure (dont l'élasticité dépend de la structure contractuelle). Les conditions d'activation sont au nombre de trois~: existence d'un modèle de revenus tiers (publicité, monétisation des données, subvention étatique), coût marginal de service proche de zéro (qui renvoie à la non-rivalité de la dimension~1), et mode d'appropriation qui supprime le signal prix marginal pour l'utilisateur final (forfait, abonnement, accès gratuit). Quand ces trois conditions tiennent, le prix au consommateur ne transmet pas le signal de rareté énergétique, et la demande agrégée de service n'est pas directement disciplinée par le coût de production.}


\paragraph{La réallocation du temps et la mutation de la borne active}

\markREG

Le deuxième canal est le temps. Quand le prix d'un service est nul et que la contrainte de budget ne mord plus, le temps disponible devient la borne active de la consommation. \textcite{goolsbee_valuing_2006} proposent de valoriser les produits numériques par le temps passé plutôt que par la dépense, ce qui donne un ordre de grandeur radicalement différent de la valeur attribuée au secteur. \textcite{kopytov_cheap_2023} montrent que la baisse du prix du loisir numérique agit comme un bien complémentaire au temps libre~: le numérique rend le loisir si attractif que le coût d'opportunité du travail change, ce qui contribue à la baisse des heures travaillées. \textcite{aguiar_leisure_2021} documentent un résultat convergent sur les jeunes hommes américains dont l'offre de travail a baissé de manière disproportionnée entre 2004 et 2015, corrélée à l'adoption des jeux vidéo et du streaming. \textcite{rachel_leisure_2024} formalise ce canal dans un modèle de croissance macroéconomique où les services numériques gratuits sont fournis en échange de l'attention~; le modèle démontre une baisse structurelle de l'offre de travail et explique la non-traduction du progrès numérique dans le PIB mesuré. Ce modèle est compatible avec un cadre DSGE et donc potentiellement traduisible en contrainte dans un modèle d'équilibre général.

Le canal du temps a un précédent historique à identification causale dans le matériau de l'analyse historique. \textcite{gentzkow_television_2006} montre que la télévision réduit la participation électorale par crowding out de l'information politique~; \textcite{stromberg_radio_2004} montre que la radio produit l'effet inverse. Le contraste entre les deux médiums~--- même modèle de financement (publicité), effets opposés sur l'allocation du temps~--- établit que le canal du temps n'est pas une conséquence mécanique du prix nul~; il dépend de la qualité de substitution du nouveau service par rapport aux activités qu'il déplace.

La combinaison des deux premiers canaux produit un effet que ni l'un ni l'autre ne produit seul. Un service payant dont le coût baisse génère un rebond classique au sens de \textcite{jevons_coal_1865}~: la consommation augmente parce que le prix diminue, mais le prix reste un signal qui borne l'expansion. Un service gratuit supprime ce signal~; le temps disponible devient la borne active. Quand en plus le service numérique est conçu pour capter l'attention~--- \textcite{allcott_welfare_2020} montrent que les utilisateurs de Facebook valorisent le service (ils exigent plus de 100~dollars pour y renoncer pendant un mois) mais que la déconnexion améliore leur bien-être subjectif~--- le temps lui-même n'est plus librement alloué. Les préférences révélées et les préférences expérimentées divergent.

Le numérique agentique (section~\ref{subsec:agentique}) introduit une mutation de ce canal. Tant que le numérique reste un service consommé par un humain~--- on regarde une vidéo, on écrit un message, on cherche une information~---, le temps humain est la borne active~: on ne regarde pas deux vidéos à la fois. L'IA agentique relâche partiellement cette borne~: une tâche humaine peut devenir un processus computationnel composé de plusieurs inférences enchaînées, et la demande de compute n'est alors plus proportionnelle au temps passé par l'utilisateur. La contrainte passe du temps de l'utilisateur au budget de calcul, de l'attention humaine à la capacité d'orchestration, de la rareté cognitive à la rareté infrastructurelle. Si ce régime se généralise, la demande de compute pourrait n'être bornée ni par le prix (H-D5.1~: prix nul), ni par le temps (H-D5.2~: borne relâchée), ce qui pose la question d'une troisième borne~--- l'infrastructure énergétique~--- dont les conditions relèvent de la dimension~9 (géographie) et de la dimension transversale 0-TER (énergie). Ce point reste conjectural~: les mesures empiriques de la demande machine-to-machine sont encore fragmentaires, et l'ampleur du phénomène n'est pas établie.

\paragraph{Hypothèse H-D5.2} \emph{La réallocation du temps vers le loisir numérique modifie les préférences et l'offre de travail par un canal absent du cadre standard, dans lequel le temps n'est pas une contrainte active de consommation. Les conditions d'activation sont la combinaison d'un prix de service nul ou faible, d'une qualité de substitution élevée par rapport aux activités déplacées, et d'un service dont l'architecture maximise le temps passé. Quand ces conditions tiennent, le temps disponible devient la borne active de la demande agrégée. L'IA agentique pourrait relâcher cette borne en déplaçant la contrainte du temps humain vers le budget de calcul~; cette mutation reste une conjecture dont l'ampleur est à établir.}


\paragraph{La demande induite par l'architecture informationnelle}

\markEVO \markINST

Le troisième canal est celui de la production de la demande par le système qui la sert. Dans le cadre standard, la demande est l'expression de préférences préexistantes. Le troisième mécanisme pose que dans les marchés informationnels, une part significative de la demande est produite par l'architecture du service~--- et que cette part a des conséquences directes en compute et en énergie.

L'analyse historique suggère que ce canal est aussi ancien que les techniques de l'information elles-mêmes. IBM, dans la période mécanographique, ne répond pas à une demande de calcul existante~: l'ingénieur Bolle au PLM reformate les processus de l'entreprise, crée le besoin de traitement mécanisé, puis y répond par l'équipement propriétaire. Les amateurs radio des années~1920 forcent le passage du point-à-point au broadcasting contre la volonté de l'industrie~: \textcite{douglas_inventing_1987} montre que Marconi, les compagnies de câbles et le Post Office britannique ne voyaient pas la valeur de la diffusion à tous. La station KDKA de Pittsburgh (novembre~1920) commence à émettre pour stimuler la vente de récepteurs~: le service crée la demande de l'équipement, pas l'inverse. L'État joue un rôle symétrique~: la Navy finance le réseau Pacifique avant qu'une demande commerciale n'existe~; le Census crée le marché de la tabulatrice avant que les entreprises ne sachent qu'elles en ont besoin.

Le phénomène contemporain qui rend ce canal central est documenté par la séquence 2001--2008. \textcite{zuboff_surveillance_2019} propose du mécanisme le terme de \emph{surplus comportemental}~: les traces d'usage (clics, requêtes, durées, abandons) sont des actifs prédictifs valorisables dont l'agrégation forme la matière première d'un nouveau type de marché. La capture comportementale exige des serveurs en fonctionnement permanent~: le lien entre architecture de la demande et consommation de compute est direct. Le lancement de l'iPhone en juin~2007 puis de l'App Store en juillet~2008 transforme chaque séquence d'interaction en enregistrement valorisable~: l'architecture du terminal devient un instrument de production de la demande.

Trois cadres théoriques éclairent des aspects différents du même phénomène. \textcite{anderson_competition_2012} modélisent l'attention des consommateurs comme une ressource commune (\emph{common-property resource}) sur-exploitée par les plateformes et les annonceurs, créant une «~congestion informationnelle~» qui dégrade le bien-être social. Le modèle est un modèle d'équilibre de marché qui traite l'attention comme un input rare~--- le parallèle avec le spectre électromagnétique, premier bien commun congestionnable de l'histoire de l'ICT (section~\ref{subsec:broadcasting}), est structurel~: le spectre est un bien commun congestionnable physique, l'attention est un bien commun congestionnable cognitif, et les deux subissent le même type de sur-exploitation. \textcite{de_corniere_model_2019} formalisent comment un algorithme de recommandation choisit de biaiser ses résultats pour maximiser ses profits plutôt que l'utilité de l'utilisateur~: l'architecture de choix agit comme un mécanisme de tarification non linéaire imposé aux fournisseurs de services. \textcite{allcott_digital_2022} quantifient causalement la part de la demande qui est produite par le système plutôt que par les préférences~: dans une expérimentation randomisée sur les réseaux sociaux, 31~\% de l'usage est attribuable à des problèmes de self-control (biais de présent, formation d'habitudes), ce qui signifie qu'un tiers de la demande~--- et du compute associé~--- ne correspond pas à l'expression de préférences stables mais à une demande artificiellement stimulée par l'architecture du service.

L'idée que les préférences sont endogènes au système économique n'est pas sans précédent. \textcite{bowles_endogenous_1998} formalise dans le \emph{Quarterly Journal of Economics} les conséquences culturelles des marchés et des institutions économiques sur les préférences des agents. \textcite{galbraith_affluent_1958} avait diagnostiqué l'\emph{effet de dépendance} par voie sociologique~: dans une économie d'abondance, la production crée les besoins qu'elle satisfait. Ce qui est propre à l'approche ici développée, c'est l'articulation entre trois éléments que ces auteurs ne connectent pas~: le mode d'appropriation du capital (H-D2.2~: données hébergées, format propriétaire, coûts de sortie asymétriques), l'architecture de l'attention (Anderson et de Palma~: congestion informationnelle), et la conséquence en compute (chaque unité de demande induite consomme de l'énergie d'inférence ou de service). La demande induite par architecture n'est pas seulement un phénomène de bien-être~; elle a un coût physique en compute et en énergie qui est d'autant plus significatif que le coût marginal d'inférence croît avec la complexité des modèles déployés.

Le mécanisme a un analogue dans l'économie des infrastructures physiques. \textcite{duranton_fundamental_2011} estiment causalement que l'élasticité de l'usage des routes par rapport à la capacité est exactement de~1~: créer 10~\% de capacité routière supplémentaire induit 10~\% de trafic supplémentaire. C'est la Loi de Say appliquée à l'infrastructure. \textcite{nordhaus_two_2007} documente un mécanisme structurellement comparable pour le calcul~: le coût du calcul a baissé d'un facteur $10^{15}$ en deux siècles, et le volume a crû dans une proportion comparable. Personne en 1900 ne «~demandait~» un trillion de FLOPS~; la demande a été produite par la disponibilité de la capacité. La question ouverte est de savoir si cette élasticité unitaire~--- attestée pour les routes, suggérée historiquement pour le compute~--- tient dans le régime contemporain de l'IA, où le coût d'inférence est un coût marginal positif et croissant. \textcite{coroama_digital_2019} montrent que les gains d'efficacité énergétique des TIC sont systématiquement compensés par des effets d'induction~: l'infrastructure crée de nouveaux usages qui n'auraient pas existé autrement. L'estimation causale à la Duranton-Turner pour le compute reste un trou dans la littérature~; la conjecture est posée ici, pas testée.

\paragraph{Hypothèse H-D5.3} \emph{Dans les marchés informationnels, une part significative de la demande est produite par l'architecture du système qui la sert, et non par l'expression de préférences préexistantes. Les conditions d'activation sont la combinaison d'un modèle de revenus tiers (condition partagée avec H-D5.1), d'une accumulation d'actif non portable côté utilisateur (condition de H-D2.2), et d'un service conçu pour maximiser l'engagement plutôt que satisfaire une préférence exprimée. Quand ces conditions sont réunies, la demande de compute est endogène au système, et les gains d'efficacité énergétique sont compensés par l'induction de nouveaux usages.} Cette hypothèse articule un résultat empirique récent~--- les 31~\% de \textcite{allcott_digital_2022}~--- avec un mécanisme historiquement documenté~--- la production de la demande par le fournisseur dans la mécanographie et le broadcasting~--- et avec un analogue formel dans l'économie des infrastructures~--- la loi fondamentale de la congestion de \textcite{duranton_fundamental_2011}. L'estimation causale de l'élasticité de la demande de compute par rapport à la capacité reste un travail empirique à réaliser.


\subsubsection{Test contemporain}

\paragraph{Pour H-D5.1 (dissociation prix-coût et asymétrie d'élasticité).} \textcite{aubouin_economie_2023} formalise le canal dans un modèle de croissance où un secteur numérique oligopolistique offre des services gratuits financés par la publicité et les données~; la qualité du service augmente la demande de biens finaux, de sorte que la concentration du secteur numérique peut être bénéfique pour le bien-être agrégé~--- un résultat conditionnel aux hypothèses du modèle, mais qui montre que la dissociation n'est pas une imperfection à corriger~: c'est un design. Sur données INSEE 2007--2019, les ménages au forfait illimité mobile ne modulent plus leur consommation avec le prix~; l'élasticité-prix de la consommation tombe. \textcite{axenbeck_what_2022} confirment côté firmes~: les entreprises plus numérisées sont moins sensibles au prix de l'énergie. \textcite{de_corniere_data_2025} montrent que les données ont deux effets opposés~--- un \emph{mark-up effect} et un \emph{surplus-extraction effect}~--- dont la résultante détermine qui supporte effectivement le coût de l'infrastructure de collecte.

La dissociation a une conséquence macroéconomique que \textcite{vanark_productivity_2016} identifie sous un autre angle~: la transition d'une économie numérique productrice (hardware, logiciels, gains de productivité mesurables) vers une économie numérique consommatrice (plateformes, applications, surplus non mesuré) modifie la composition de la demande finale sans que la comptabilité nationale ne le capte. \textcite{byrne_digital_2022} et \textcite{syverson_challenges_2017} montrent que le \emph{mismeasurement} n'explique qu'une fraction marginale de ce décalage~: l'essentiel est structurel. L'inférence IA accentue la dissociation~: le coût marginal d'une requête ChatGPT est environ dix fois celui d'une recherche Google classique (section~\ref{subsec:genai}), mais le prix pour l'utilisateur du tier gratuit reste nul.

\paragraph{Pour H-D5.2 (réallocation du temps et mutation de la borne).} Les résultats de \textcite{kopytov_cheap_2023} et \textcite{aguiar_leisure_2021} confirment le canal du temps dans le régime pré-agentique. L'IA agentique introduit un changement de régime potentiel~: quand un agent IA enchaîne plusieurs inférences pour accomplir une tâche humaine, la demande de compute est multipliée sans que le temps humain n'augmente. La section~\ref{subsec:agentique} documente le déplacement~: la contrainte passe du temps de l'utilisateur au budget de calcul. Les mesures empiriques de ce déplacement sont encore fragmentaires, et les projections d'Epoch~AI et de l'IEA sur la demande d'inférence reposent sur des hypothèses de diffusion incertaines~; le point est posé comme tendance structurelle, pas comme fait établi.

\paragraph{Pour H-D5.3 (demande induite par architecture).} \textcite{allcott_digital_2022} fournissent la mesure causale la plus directe~: 31~\% de l'usage des réseaux sociaux est attribuable à des problèmes de self-control. \textcite{anderson_competition_2012} donnent le cadre formel~: l'attention comme ressource commune sur-exploitée. Le cas contemporain le plus prononcé est celui des plateformes vidéo~: l'autoplay, les recommandations algorithmiques et le design sans friction sont des mécanismes de production de la demande dont l'effet sur le compute est direct (chaque vidéo regardée consomme de la bande passante, du stockage et du compute de recommandation). L'ancien parallèle historique le plus direct est la station KDKA qui émet pour stimuler la vente de récepteurs~: le service crée la demande de l'équipement. Le test négatif est celui des communs numériques (D2, H-D2.2)~: quand le mode d'appropriation est ouvert (logiciel libre, format standard, données portables), la condition d'activation n'est pas remplie et la demande induite est moindre~--- l'utilisateur n'accumule pas d'actif non portable, et le service n'est pas conçu pour maximiser l'engagement.


\subsubsection{Fait stylisé}

L'ICT transforme la structure de la demande finale par trois canaux distincts~: la dissociation entre le prix perçu et le coût de production (H-D5.1), qui rompt le bouclage prix-rareté et crée une asymétrie d'élasticité~; la réallocation du temps vers le loisir numérique (H-D5.2), qui modifie les préférences et l'offre de travail par un canal que les élasticités-prix ne captent pas~; et la production de la demande par l'architecture du système (H-D5.3), qui rend la demande de compute endogène au service qui la sert. Les trois canaux opèrent conjointement~: le service gratuit (H-D5.1) libère du temps (H-D5.2), le temps libéré est capturé par une architecture conçue pour maximiser l'engagement (H-D5.3), et la demande résultante n'est bornée ni par le prix ni par les préférences. Ce fait stylisé croise les dimensions~1 (la non-rivalité rend le prix nul possible), 2 (le mode d'appropriation crée le lock-in côté demande), 3 (l'offre de travail est affectée par la réallocation du temps) et~6 (le signal de prix est structurellement supprimé).

Les trois mécanismes exigent qu'un modèle énergie-économie représente explicitement la consommation à prix nul et différencie l'élasticité-prix selon le type d'acteur~--- faible pour les utilisateurs finaux numérisés au sens d'Aubouin, significative pour les opérateurs d'infrastructure. H-D5.1 et l'asymétrie d'élasticité peuvent être traités comme paramétrage de scénario via la calibration sectorielle des élasticités, sous réserve d'examen dans la modélisation. Le canal de la réallocation du temps (H-D5.2) reste comme contrainte interprétative~; la mutation agentique du canal est une conjecture qui informera la construction des scénarios sans s'y traduire mécaniquement tant que les données empiriques manquent. La demande induite (H-D5.3) peut être approchée par la calibration de scénarios différenciant un régime à demande exogène (préférences stables, élasticités calibrées) et un régime à demande endogène (élasticité unitaire de la demande de compute par rapport à la capacité, à la Duranton-Turner)~; la différence entre les deux scénarios est elle-même un résultat analytique.


\begin{table}[!ht]
\caption{Les trois mécanismes de la dimension et les hypothèses associées.}\label{tab:d5-mecanismes}
\centering\footnotesize
\renewcommand{\arraystretch}{1.25}
\begin{tabularx}{\textwidth}{@{}p{2.6cm}p{2.8cm}p{3.6cm}X@{}}
\toprule
\textbf{Mécanisme} & \textbf{Canal mobilisé} & \textbf{Conditions d'activation} & \textbf{Hypothèse} \\
\midrule
M1. Dissociation prix-coût & Bouclage prix-rareté & Modèle de revenus tiers \emph{et} coût marginal proche de zéro \emph{et} mode d'appropriation supprimant le signal prix marginal & H-D5.1~: le prix ne transmet pas le signal de rareté énergétique~; asymétrie d'élasticité entre utilisateurs finaux et opérateurs \\
M2. Réallocation du temps & Contrainte active & Prix nul \emph{et} qualité de substitution élevée \emph{et} architecture maximisant le temps passé & H-D5.2~: le temps devient la borne active~; l'IA agentique pourrait déplacer cette borne vers le budget de calcul \\
M3. Demande induite par architecture & Nature des préférences & Modèle de revenus tiers \emph{et} actif non portable côté utilisateur \emph{et} service conçu pour l'engagement & H-D5.3~: la demande de compute est endogène au système~; les gains d'efficacité sont compensés par l'induction de nouveaux usages \\
\bottomrule
\end{tabularx}
\renewcommand{\arraystretch}{1.0}
\end{table}

Une interaction entre mécanismes mérite d'être signalée. H-D5.3 (demande induite) opère par le même canal que H-D2.2 (mode d'appropriation) mais du côté de la demande plutôt que de l'offre~: le design du capital qui crée le lock-in côté firme crée aussi la demande captive côté utilisateur. Cette interaction D2-D5 est la boucle de rétroaction la plus directe entre les dimensions offre et demande~; elle sera développée dans la section d'intégration.


% =============================================================================
\clearpage

% --- FIN DE LA SECTION D5 ---
